% VI. Експериментални резултати и анализ.
% ВНИМАНИЕ: този файл съдържа САМО формата на резултатите. Нито едно число тук не се
% пише „приблизително“ или „за пример“, попълва се от изхода на измерванията, иначе
% числото оцелява до предаването и никой не помни, че е било измислено.
\section{Експериментални резултати и анализ}
\label{sec:results}

Разделът представя резултатите, получени по методиката от Раздел~\ref{sec:method}.
Всяка таблица и всяка фигура се придружава с коментар в текста и с препратка по А.8 и
А.9. Резултат без тълкуване не се включва.

\subsection{Проверка на условията за валидност}
\label{subsec:res-validity}

Преди същинските резултати се проверява, че условията от §5.2 са изпълнени: машината
не изпълнява друго натоварване, разсейването между повторенията е в очакваните
граници и резултатите при една нишка са възпроизводими при фиксирано начално число.

\TODO{Резултат от проверката. Ако някое условие не е изпълнено, това се записва тук,
а не се премълчава.}

\subsection{Ускорение и ефективност спрямо броя нишки}
\label{subsec:res-speedup}

\TODO{Таблица: редове по брой нишки (1, 2, 4, 8, ... до броя логически ядра), колони
по схема за паралелно изпълнение. Клетка: медианно време и междуквартилен размах.}

\TODO{Таблица: същата структура, но с ускорение и ефективност, изчислени спрямо
изпълнението с една нишка.}

\TODO{Фигура: ускорение спрямо брой нишки, с нанесена линия на идеалното ускорение
като отправна линия. Легендата е задължителна по А.9.}

\TODO{Коментар: къде кривата се отклонява от идеалната и коя измерена величина
обяснява отклонението. Обяснение без подкрепящо измерване не се пише.}

\subsection{Качество на решението при фиксиран бюджет}
\label{subsec:res-quality}

\TODO{Таблица: редове по инстанция, колони по метаевристика. Клетка: медианно
относително отклонение от известната стойност.}

\TODO{Фигура: разпределение на постигнатото качество по повторения, за да се вижда
разсейването, а не само медианата.}

\TODO{Коментар: подрежда ли се качеството еднакво при всички инстанции, или
подредбата зависи от размера на инстанцията.}

\subsection{Сравнение между паралелните схеми}
\label{subsec:res-cooperation}

\TODO{Таблица: за всяка схема, постигнато качество при равен изчислителен бюджет и
измерено време, прекарано в синхронизация.}

\TODO{Коментар: изплаща ли се комуникацията. Отговорът е валиден само за измерените
инстанции и се формулира с това ограничение.}

\subsection{Разпределение на времето и профил}
\label{subsec:res-profile}

\TODO{Таблица: дял на времето по участъци, тоест оценяване на целевата функция,
генериране на съседи, синхронизация, останало.}

\TODO{Съпоставка на измерената производителност на горещия участък с горната граница
по модела на покривната линия и извод дали участъкът е ограничен от паметта или от
изчисленията.}

\subsection{Обобщение на резултатите}
\label{subsec:res-summary}

\TODO{Отговор на всеки от четирите въпроса от §5.1, по един абзац. Въпрос, на който
измерванията не дават отговор, се отбелязва като неотговорен, а не се заобикаля.}
